%% kapa: exam=1 round=28 year=2017 period=2 form=A subject=accounting subject_ko=회계학 seq=21 number=61
%% source: Q-Net 감정평가사 기출문제 (cst00309) artlSeq=5199603 file=extracted/1008/5199603/2017년 제28회 감정평가사 1차시험 문제지/2교시 A형.pdf
%% conversion: kordoc 4.24.51 -> tools/gen1.py (auto); review=proofread-sonnet; flags=none
\begin{question}{61}
유용한 재무정보의 질적 특성에 관한 설명으로 옳지 않은 것은?
\begin{choices}
  \item 재무정보가 유용하기 위해서는 목적적합해야 하고 나타내고자 하는 바를 충실하게 표현해야 한다.
  \item 보강적 질적 특성을 적용하는 것은 어떤 규정된 순서를 따르지 않는 반복적인 과정이므로 때로는 하나의 보강적 질적 특성이 다른 질적 특성의 극대화를 위해 감소되어야 할 수도 있다.
  \item 회계기준위원회는 중요성에 대한 획일적인 계량 임계치를 정하거나 특정한 상황에서 무엇이 중요한 것인지를 미리 결정할 수 있다.
  \item 중요성은 개별 기업 재무보고서 관점에서 해당 정보와 관련된 항목의 성격이나 규모 또는 이 둘 모두에 근거하여 해당 기업에 특유한 측면의 목적적합성을 의미한다.
  \item 근본적 질적 특성을 충족하면 어느 정도의 비교가능성은 달성될 수 있을 것이다.
\end{choices}
\end{question}
\answer{3}
